
\section{Giry Monad}

\DPTtodo[inline]{Are we deleting this?}
\subsection{Category theory}
A category $\mathcal{C}$ is a collection of objects and morphisms between them. So if $A,B,C$ and $D$ are objects in $\mathcal{C}$, then $f : A \to B$ and $g : B \to C$ and $h : C \to D$ are examples of morphisms. They satisfy the associativity rule $f \circ (g \circ h) = (f \circ g) \circ h$ and for every object $B$ there exist an identity morphism $\operatorname{id}_B : B \to B$ with the property $f \circ \operatorname{id}_A = f$ for all $f$ and $\operatorname{id}_B \circ g = g$ for all $g$. A \emph{functor} is a mapping between categories. The functor $F : \mathcal{C} \to \mathcal{D}$ takes every object in $\mathcal{C}$ to an object in $\mathcal{D}$ so that the structure of morphism composition is preserved. An \emph{endofunctor} is a functor from a category to itself, and a \emph{bifunctor} is a functor from a pair of categories to another category, i.e. $F : \mathcal{A} \times \mathcal{B} \to \mathcal{C}$. A \emph{natural transformation } is a mapping between functors. More precisely, for two functors $F,G : \mathcal{C} \to \mathcal{D}$, a natural transformation $\epsilon : F \to G$ associates to every object $C$ in $\mathcal{C}$ a morphism $\epsilon_C : F(C) \to G(C)$ in $\mathcal{D}$. A \emph{monoidal category} $\mathcal{C}$ is a category with some additional monoidal structure, namely, an identity object $I$ and a bifunctor $\otimes : \mathcal{C} \times \mathcal{C} \to \mathcal{C}$ called the \emph{tensor product}, together with other natural isomorphisms providing associativity of tensor products, right and left identity with the identity object $I$. A \emph{monoid} $(M,\mu,\eta)$ in a monoidal category $\mathcal{C}$ is an object $M$ in $\mathcal{C}$ together with two morphisms (with some properties): the associative binary operator $\mu : M \otimes M \to M$ and the identity $\eta : I \to M$. A \emph{monad} is a monoid in the category of endofunctors. This is, consider the category $\mathcal{F}$ whose objects are endofunctors and whose morphisms are natural transformations between them. This is a monoidal category: there exist an identity endofunctor $1_{\mathcal{F}}(F)=F$ for all $F \in \mathcal{F}$ plus a tensor product $\otimes : \mathcal{F} \times \mathcal{F} \to \mathcal{F}$ defined by functor composition such that the required associativity and identity properties hold. $\mathcal{F}$ is therefore a monoidal category and any specific monoid $(F,\mu,\eta)$ is a monad.
\subsection{Probability}
A \emph{measurable space} $(X,\mathcal{A})$ is a set $X$ equipped with a $\sigma$-algebra $\mathcal{A}$ (consisting of measurable subsets of $X$). A measure $\mu : \mathcal{A} \to \mathbb{R}_{0\geq}$ is a set function from the $\sigma$-algebra to non-negative reals. If $\mu(X)=1$, we say $\mu$ is a \emph{probability measure}. A measurable space together with a probability measure $(X,\mathcal{A}, \mu)$ is called a \emph{probability space}.
A measurable map $T : (X,\mathcal{A}) \to (Y,\mathcal{B})$ is a map $T : X \to Y$ so that for any $B \in \mathcal{B}$, $T^{-1}(B)$ is in $\mathcal{A}$.
Let $\operatorname{Meas}$ be the category whose objects are all measurable spaces and morphisms are measurable mappings between them. For a fixed measurable space $M$ in $\operatorname{Meas}$, we can consider the space of all possible probability measures that could be placed on it and denote it by $\mathcal{M}_1(M)$ (so $\mathcal{M}_1(M)$ is a space of probability measures).
The point is that $\mathcal{M}_1(M)$ itself is also a measurable space. \\
For a measurable space $(X,\mathcal{A})$, and any $A \in \mathcal{A}$, we can define $\tau_A(\nu)=\nu(A)=\int_{X} \chi_{A}(x) d\nu(x)$ where $\chi_A$ denotes the characteristic function of $A$. Now, consider $T : (X, \mathcal{A}) \to (\mathbb{R},\mathcal{B}(\mathbb{R}))$ where $\mathcal{B}(\mathbb{R})$ is the Borel-$\sigma$-algebra for the usual Euclidean topology of $\mathbb{R}$). Define $\tau_T : \mathcal{M}_1(M) \to \mathbb{R}$ as $\tau_T(\nu)=\int_X T(x) d\nu(x)$. We call these maps $\tau_A$ and $\tau_T$ the \emph{evaluation maps}.
For a measurable space $M$, there exist a topology on $\mathcal{M}_1(M)$ called the weak$^*$-topology that makes all evaluation maps (for any measurable set $A$ and any measurable map $T$) continuous. From there, we can generate the Borel $\sigma$-algebra $\mathcal{B}(\mathcal{M}_1(M))$ that makes the evaluation functions $\tau_A$ and $\tau_T$ measurable. Then, $(\mathcal{M}_1(M),\mathcal{B}(\mathcal{M}_1(M))$ is itself a measurable space and thus an object in $\operatorname{Meas}$. The space $\mathcal{M}_1(M)$ has properties that carry over from $M$ (convexity, completeness, compactness, etc. Maybe write about this? Millman and Choquet theorems of functional analysis here?). \\ \\

Now, let $T : M \to N$ be a measurable mapping between two measurable spaces $M$ and $N$ (i.e. a morphism in $\operatorname{Meas}$). For any $\nu \in \mathcal{M}_1(M)$ we can define $\mathcal{T}(\nu)\coloneqq\nu \circ T^{-1}$ (pushforward measure of $\nu$ under $T$), an element in $\mathcal{M}_1(N)$.
Therefore, $\mathcal{M}_1$ is an endofunctor of $\operatorname{Meas}$.
In fact, we shall show that $\mathcal{M}_1$ is a monad. 

\begin{enumerate}
\item Identity: Let $M$ be a measurable space in $\operatorname{Meas}$. Define $\eta(x)(A)\coloneqq\chi_A(x)$ for any $x \in M$ and any measurable set $A \subset M$. This defines a probability measure on $M$ assigning 1 to measurable sets that contain $x$ and $0$ to those that don't. This yields $\eta : M \to \mathcal{M}_1(M)$. Note that $\eta$ takes points in measurable sets to probability measures on those spaces (points to their corresponding Dirac measures).
\item Multiplication: recall that the category of endofunctors $\mathcal{F}$ is monoidal, so there exists a tensor product $\otimes : \mathcal{F} \times \mathcal{F} \to \mathcal{F}$. Take $M$ a measurable space in $\operatorname{Meas}$ and consider $\mathcal{M}_1(M)$. Consider $\mathcal{M}_1(\mathcal{M}_1(M))$. Since $\mathcal{M}_1$ is an endofunctor, this is again a measurable space, and for any measurable subset $A$ of $M$ we have a family of maps $\tau_A$ taking a probability measure $\mathcal{M}_1(\mathcal{M}_1(M))$ and evaluating it on $A$. 
Take $\rho\in \mathcal{M}_1(\mathcal{M}_1(M))$ and let $A$ be a measurable subset of $M$. Define $\mu$ as 
\[
  \mu(\rho)(A)\coloneqq\int_{\mathcal{M}_1(M)} \tau_A d\rho.
\]
So, $\mu(\rho) : \mathcal{A} \to \mathbb{R}$ is a probability measure on $M$.

Thence, $\mu_M : \mathcal{M}_1(\mathcal{M}_1(M)) \to \mathcal{M}_1(M)$.

We can define, similarly,
\[
  \mu(\rho)(A)\coloneqq\int_{\mathcal{M}_1(M)} f_A d\rho.
\]

Thus, we get a natural transformation $\mu : \mathcal{M}_1^2 \to \mathcal{M}_1$.

\end{enumerate}


%%% Local Variables:
%%% mode: latex
%%% TeX-master: "Smoothings"
%%% End:
